% Standalone problem document, generated from the adjacent README.md.
% Compile with XeLaTeX. Mathematical content is maintained in Markdown.
\documentclass[11pt,a4paper]{article}
\usepackage[fontsize=10.5pt]{scrextend}
\usepackage[margin=22mm,headheight=15pt,headsep=9mm,footskip=11mm]{geometry}
\usepackage{fontspec}
\setmainfont{texgyrepagella}[Extension=.otf,UprightFont=*-regular,BoldFont=*-bold,ItalicFont=*-italic,BoldItalicFont=*-bolditalic]
\setsansfont{texgyreheros}[Extension=.otf,UprightFont=*-regular,BoldFont=*-bold,ItalicFont=*-italic,BoldItalicFont=*-bolditalic]
\setmonofont{texgyrecursor}[Extension=.otf,UprightFont=*-regular,BoldFont=*-bold,ItalicFont=*-italic,BoldItalicFont=*-bolditalic,Scale=MatchLowercase]
\usepackage{amsmath,amssymb}
\usepackage{unicode-math}
\setmathfont{texgyrepagella-math.otf}
\usepackage{xcolor,microtype,enumitem,needspace,fancyhdr}
\usepackage{bookmark}
\definecolor{ink}{HTML}{17344B}
\definecolor{accent}{HTML}{197D80}
\definecolor{muted}{HTML}{596773}
\hypersetup{colorlinks=true,linkcolor=accent,urlcolor=accent,pdftitle={NM-03 - Complexity
of globally optimal nonnegative rank-two
approximation},pdfauthor={Open Problems in Numerical Linear Algebra}}
\urlstyle{same}
\setlength{\parindent}{0pt}
\setlength{\parskip}{5pt plus 1pt minus 1pt}
\linespread{1.04}
\setlength{\emergencystretch}{3em}
\widowpenalty=10000
\clubpenalty=10000
\displaywidowpenalty=10000
\allowdisplaybreaks[1]
\setlist{nosep,leftmargin=1.5em}
\providecommand{\tightlist}{\setlength{\itemsep}{0pt}\setlength{\parskip}{0pt}}
\providecommand{\pandocbounded}[1]{#1}
\pagestyle{fancy}
\fancyhf{}
\fancyhead[L]{\sffamily\footnotesize\color{muted}OPEN PROBLEMS IN NUMERICAL LINEAR ALGEBRA}
\fancyhead[R]{\sffamily\bfseries\footnotesize\color{accent}NM-03}
\fancyfoot[L]{\parbox[b]{0.9\textwidth}{\sffamily\fontsize{8}{10}\selectfont\color{muted}Editorial ratings; a bounded search cannot certify that a problem remains unsolved.\\Literature check: 2026-09-08}}
\fancyfoot[R]{\sffamily\footnotesize\color{muted}\thepage}
\renewcommand{\headrulewidth}{0.3pt}
\renewcommand{\headrule}{\hbox to\headwidth{\color{accent}\leaders\hrule height \headrulewidth\hfill}}
\makeatletter
\renewcommand{\section}{\Needspace{5\baselineskip}\@startsection{section}{1}{0pt}{12pt plus 2pt minus 2pt}{4pt}{\sffamily\bfseries\large\color{ink}}}
\renewcommand{\subsection}{\Needspace{4\baselineskip}\@startsection{subsection}{2}{0pt}{9pt plus 1pt minus 1pt}{3pt}{\sffamily\bfseries\normalsize\color{ink}}}
\makeatother
\setcounter{secnumdepth}{0}
\begin{document}
{\sffamily\small\color{muted}Nonnegative and positive
factorizations\par}
\vspace{3pt}
{\raggedright\hyphenpenalty=10000\exhyphenpenalty=10000\sffamily\bfseries\fontsize{21}{24}\selectfont\color{ink}Complexity
of globally optimal nonnegative rank-two approximation\par}
\vspace{7pt}
{\sffamily\small\textbf{Difficulty:} extreme\quad\textbf{Importance:} broadly
interesting\par}
\vspace{4pt}
{\color{accent}\hrule height 0.8pt}
\vspace{6pt}
\textbf{Topic:} low-rank approximation; computational complexity

\textbf{Status:} open; admitted after a bounded literature search found
no resolution.

\subsection{Problem statement}\label{problem-statement}

Given an arbitrary \(X\in\mathbb Q_+^{m\times n}\) and
\(\tau\in\mathbb Q_{\geq0}\), determine the complexity of deciding
whether

\[
\exists W\in\mathbb R_+^{m\times2},\ H\in\mathbb R_+^{2\times n}:
\quad\sum_{i=1}^m\sum_{j=1}^n
\left(X_{ij}-\sum_{k=1}^2W_{ik}H_{kj}\right)^2\leq\tau.
\]

In particular, is there a deterministic algorithm polynomial in the
total binary input length, or is this decision problem NP-hard under
polynomial-time many-one reductions? This is a complexity-classification
question, without an assumption that these two outcomes exhaust the
possibilities. The factors may have real entries; only the data and
threshold must be rational. Their inner dimension is at most two, with a
zero factor column permitted.

This fixes an exact decision interpretation of the literature's global
rank-two NMF optimization question. There is no promise that \(X\)
itself has rank two. When a rank-two truncated SVD of \(X\) is
nonnegative, a best nonnegative rank-two approximation is obtainable
from it. Arbitrary input can fall outside this tractable special case,
and alternating nonnegative least squares need not find a global
optimum.

\subsection{References}\label{references}

Gillis,
\href{https://doi.org/10.1137/1.9781611976410.ch6}{\emph{Nonnegative
Matrix Factorization}}, §6.1.3, p.~199, Theorem 6.6 and final paragraph.
Lindy, Noferini, and Van Dooren,
\href{https://arxiv.org/abs/2507.20612v1}{\emph{On rank-2 Nonnegative
Matrix Factorizations and their variants}}, §1, with §3's suboptimal
approximation and §4's ANLS initialization.

\subsection{Status check}\label{status-check}

Searches for
\textit{rank two nonnegative matrix factorization approximation NP hard polynomial time 2025 2026}
and \textit{rank-2 NMF complexity 2026} found no resolution. The July
2025 primary paper explicitly identifies rank two as an unresolved
complexity case; its contribution is an effective initial approximation
and heuristic refinement. The March 2026
\href{https://arxiv.org/abs/2603.19976}{constrained nonnegative
Gram-feasibility preprint} concerns partially specified symmetric
matrices with affine side constraints, which are absent from the problem
above.
\end{document}
